\section{Алгоритм и его реализация}\label{sec:algo}

Опишем концептуально вычислительный конвейер; технические детали и код --- в
пакетах \texttt{voronoi4d} (эталонная реализация на Python) и \texttt{combigeo}
(быстрое ядро на C++), включённых в репозиторий проекта.

\paragraph{(1) Приведение базиса.} Базис решётки предварительно
\emph{LLL-приводится} --- заменяется на эквивалентный, но <<почти
ортогональный>>, с короткими векторами (алгоритм Ленстры--Ленстры--Ловаса
\cite{LLL}). Это резко сужает все последующие переборы.

\paragraph{(2) Построение ячейки Вороного.} Грани ячейки задаются
\emph{релевантными векторами} решётки. По теореме Вороного вектор $v$
релевантен тогда и только тогда, когда $\pm v$ --- \emph{единственная} пара
кратчайших представителей класса смежности $v+2\Lambda$; отсюда релевантных
векторов не более $2(2^n-1)$. Мы находим кратчайшие представители точным
перебором ближайших векторов (\emph{CVP}, задача о ближайшем векторе, методом
сферического декодирования Финке--Поста \cite{FinckePohst}) в каждом из
$2^n-1$ ненулевых классов, проверяя попутно единственность пары. По найденным
граням перечисляются вершины ячейки. Корректность построенной ячейки
контролируется независимо: проверяется соотношение Эйлера для её $f$-вектора,
центральная симметрия множества вершин и равенство её объёма определителю
решётки. Для рассматриваемых решёток общего положения ячейка \emph{проста}: в
каждой вершине активно ровно $n$ фасет, и набор $v_1,\dots,v_n$ ниже определён
однозначно. Полнота самого списка вершин контролируется тождеством
\begin{equation}\label{eq:vertsum}
\sum_{x\in\mathrm{vert}\,V_0}\bigl|\det(v_1,\dots,v_n)\bigr|=(n+1)!,
\end{equation}
где $v_1,\dots,v_n$ --- релевантные векторы фасет, активных в вершине~$x$: слева
стоит сумма нормированных объёмов симплексов Делоне, инцидентных узлу решётки, а
она не зависит от комбинаторного типа ячейки. Для примитивной решётки все
слагаемые равны~$1$ и вершин ровно $(n+1)!$; в общем случае вершин меньше, и
\eqref{eq:vertsum} даёт проверяемое равенство вместо ссылки на связность графа.

\paragraph{(3) Расстояние точка${}\to{}$ячейка.} По
формуле~\eqref{eq:mid} нужно $\dist(v/2,V_0)$. Оно вычисляется \emph{тремя}
независимыми способами (для взаимного контроля): алгоритмом
Гилберта--Джонсона--Кирти (GJK) --- стандартным методом вычислительной
геометрии для расстояния до выпуклой оболочки \cite{GJK}, с апостериорной
оценкой погрешности через зазор двойственности; каскадом ортогональных проекций
по граням; и решением задачи квадратичного программирования по
полупространствам. На всех тестах три способа совпадают до девяти значащих цифр.

\paragraph{(4) Перебор подрешёток данного индекса.} Все подрешётки индекса $k$
взаимно однозначно соответствуют целочисленным матрицам перехода в
\emph{эрмитовой нормальной форме} (верхнетреугольным, с фиксированными
диапазонами элементов); их число равно $\sum_{d_1\cdots d_n=k}d_2d_3^2\cdots
d_n^{\,n-1}$. Перебор распараллелен по потокам (страйд по матрицам, атомарное
разделяемое текущее наилучшее значение); на 8 ядрах ускорение
${\approx}\,2{,}5\times$.

\paragraph{(5) Поиск лучшей решётки: оптимизация метрики.} Для \emph{фиксированного}
числа цветов $k$ мы не ограничиваемся классическими решётками, а
\emph{оптимизируем саму метрику}. Решётка с точностью до изометрии задаётся
своей матрицей Грама (положительно определённой квадратичной формой $Q$);
пространство таких форм разбито на конечное число \emph{вторичных конусов}, в
каждом из которых комбинаторный тип ячейки Вороного постоянен (теория приведения
Вороного; вычислительный аппарат --- Шюрман, Валлентин \cite{SchurmannVallentin};
классификация в $\R^4$ --- Валлентин, Вайсбах, Циммерман \cite{VWZ}). Мы
максимизируем $d(Q,k)$ по $Q$ безградиентным методом Нелдера--Мида \cite{NelderMead}
(в параметризации Холецкого), запуская его из многих начальных точек на пуле
процессов. Именно этот шаг --- выход в решётки \emph{общего положения} --- и
даёт улучшение оценки.

\paragraph{(6) Точная сертификация.} Численный оптимум рационализируется (форма
$Q$ с небольшими знаменателями), после чего ключевое неравенство $d>1$
проверяется в \emph{точной арифметике дробей} --- без всякой плавающей точки
(\S\ref{sec:comp}). Это переводит результат из <<вычислительного свидетельства>>
в компьютерно проверяемое доказательство.

\paragraph{Выбор оптимизатора.} Функция $d(Q,k)$ негладкая
(кусочно-постоянна по дискретной подрешётке) и многоэкстремальна, поэтому
применимы лишь безградиентные глобальные методы. Мы использовали метод
Нелдера--Мида \cite{NelderMead} с многократными рестартами, глобальный
Q-поиск (метод многократного сжатия бруса с перезапусками, \cite{GornovQ}) и
эволюционную стратегию CMA-ES \cite{Hansen}. В размерностях $n\le4$ шаг~(4)
(перебор всех подрешёток индекса $k$) дёшев, и оптимизация формы практична;
именно так получены результаты \S\S\ref{sec:main}--\ref{sec:stair}.

\subsection{Масштабируемый поиск в размерностях 5--9}\label{sec:big}

При $n\ge5$ узким местом становится шаг~(4): число подрешёток индекса $k$ растёт
как $k^{\,n-1}$ (в $\R^5$ при $k\approx140$ это $\sim10^9$), и исчерпывающая
оценка $d(Q,k)$ занимает десятки минут. Мы разработали подход без полного
перебора, основанный на переформулировке задачи как \emph{задачи о выполнимости
ограничений} (CSP).

\paragraph{Подрешётка как гомоморфизм.} Раскраска индекса $k$ задаётся
\emph{сюръективным} гомоморфизмом $\varphi\colon\Lambda\twoheadrightarrow G$ на
конечную абелеву группу $G$ порядка $k$; подрешётка $\Lambda'=\ker\varphi$
(именно сюръективность гарантирует $[\Lambda:\ker\varphi]=|G|=k$). Пусть
$F_\ell=\{v\in\Lambda\setminus\{0\}:D(v)<\ell\,\diam V_0\}$ --- конечное
<<запрещённое множество>> (при $D(v)\ge\|v\|-\diam V_0$ оно ограничено
$\|v\|<(\ell+1)\diam V_0$). Тогда
\[
d(\Lambda,\Lambda')\ge\ell\quad\Longleftrightarrow\quad
\varphi(v)\neq0\ \ \forall v\in F_\ell,
\]
т.е. нужно найти $\varphi$, <<отделяющий>> все запрещённые векторы от нуля ---
это CSP над коэффициентами форм $\varphi$.

\paragraph{Структура факторгруппы.} Структура $\Lambda/\Gamma$ определяется
выбором подрешётки, а не самой решёткой, и заранее ограничивать её нельзя.
Эмпирически у \emph{симметричных} решёток пригодные подрешётки имеют
\emph{нециклическую} факторгруппу $G=\Z/e_1\times\cdots\times\Z/e_m$,
отражающую симметрию (для $D_4$ при $k=49$ это $\Z/7\times\Z/7$; для $A_5^*$
при $k=140$ --- $\Z/2\times\Z/70$; для $E_6^*$ при $k=343$ ---
$\Z/7\times\Z/7\times\Z/7$), тогда как у найденных нами решёток общего
положения пригодной оказывалась циклическая $G=\Z/k$. Поэтому перебирать надо
\emph{все} структуры инвариантных факторов индекса $k$, а не только
циклическую.

\paragraph{Локальный поиск.} Пригодные $\varphi$ исчезающе редки (для $A_5^*/140$ их
не находит и $2{,}4\cdot10^6$ случайных проб), так как запрещённые векторы сильно
скоррелированы. Мы применяем \emph{локальный поиск} min-conflicts (как для
трудных CSP типа раскраски графов): стартуем со случайного $\varphi$ и повторно
меняем коэффициент формы, минимизируя число <<убитых>> запрещённых векторов. Это
находит нециклические подрешётки $A_5^*/140$ и $E_6^*/343$ за \emph{секунды} ---
стоимость одной оценки формы падает с десятков минут до секунд. В размерности~6
перечисление вершин ячейки ($\binom{126}{6}$) обходится расчётом расстояний
через опорные полупространства (квадратичное программирование).

\paragraph{Радиус покрытия без перечисления вершин.} Нормировка $\diam V_0=2R$
($R$ --- радиус покрытия) при $n\ge6$ тоже упирается в перечисление вершин
ячейки. На поисковой стадии мы \emph{оцениваем} $R=\max_{x\in V_0}\|x\|$ без
полного перечисления вершин: для набора случайных направлений $u$ решаются
\emph{линейные} задачи $\max\,u\!\cdot\!x$, $x\in V_0$, по
полупространственному описанию ячейки, и максимальная норма найденных опорных
точек даёт численную \emph{нижнюю} оценку $R$ --- статус \textbf{[Ч]}. Эта
оценка используется только для ранжирования кандидатов; в сертификатах теорем
радиус покрытия проверяется точно по полной структуре ячейки Вороного. Метод
безразмерен, и на тестовых решётках до $\R^9$ оценка совпала с известными
замкнутыми формами $\diam/\lambda_1$
(ср.~\cite{ConwaySloane,SchurmannVallentin}: $A_n^*{:}\ \sqrt{(n{+}2)/3}$;
$D_n{:}\ \sqrt{n/2}$; $E_6^*,E_8{:}\ \sqrt2$; $E_7{:}\ \sqrt3$) с точностью до
$5$ знаков.

\paragraph{Результаты и реализация.} Ядро метода перенесено на
C\raise.1ex\hbox{\small++} (модуль \texttt{combigeo}: \texttt{min\_conflicts},
построение $F_\ell$ через опорные полупространства, радиус покрытия); он
воспроизводит нециклическую подрешётку $E_6^*/343$ ($G=\Z/7\times\Z/7\times\Z/7$)
за ${\approx}0{,}3$~с. Метод валидирован: он воспроизводит известные
конструкции $A_5^*/140$ и $E_6^*/343$ (нецикличные, с точной проверкой структуры
факторгруппы), причём в $\R^5$ поиск \emph{минимального} индекса на решётке
$A_5^*$ находит именно $140$ как наименьший пригодный индекс (карта пригодных
индексов при $\ell=1$: $140,144,156,160,164,166,168,172,\dots$, далее плотно с
$183$), а при $k<140$ пригодных подрешёток $A_5^*$ не найдено. Это согласуется с
предложением~12 из \cite{ABPR}, где полным перебором эрмитовых нормальных форм
доказано $\min_{\Lambda'}\chi_s(\E^5,A_5^*,\Lambda',2R)=140$; отметим, что там
минимум берётся по подрешёткам \emph{фиксированной} решётки $A_5^*$ и лишь при
$\ell=2R$, тогда как наш поиск ведётся при $\ell=1$. Для $E_6^*$ аналогичного
утверждения о минимальности в \cite{ABPR} нет: оценка $343$ получена там
конструктивно (теорема~3), без перебора.


\paragraph{Почему прежний барьер $140$ оказался артефактом цели.}
Первый вариант оптимизации по формам Грама использовал бинарную цель ---
минимальное число неотделённых векторов $\mathrm{killed}(B,k)$. При
$k=139$ он неизменно оставлял один конфликт, поэтому прежняя кампания ошибочно
указывала на устойчивость порога~$140$. Бинарная цель не различает глубокий и
почти устранённый конфликт и, главное, после деформации формы не возвращает
новую геометрическую информацию дискретному оракулу.

Новый поиск сочетает четыре механизма. Во-первых, целочисленный кандидат
доводится до требуемого определителя точной последней поправкой
\[
\det(A+\delta E_{ij})=\det A+\delta\,\operatorname{cof}_{ij}(A).
\]
Во-вторых, при фиксированном ядре непосредственно максимизируется
$\min D/\diam V$ в параметризации
$B=B_0\exp H$, $H=H^{\mathsf T}$, $\operatorname{tr}H=0$, с температурным продолжением сглаженного минимума. В-третьих, почти активные проекции и дальние вершины
образуют локальную max--min ЛП в доверительной области, а каждый принятый шаг повторно
проверяется полным геометрическим оракулом. Наконец, после локального потолка
ядро \emph{снова} оптимизируется на уже деформированной метрике.

Для простого индекса $p$ ядро задаётся строкой
$a\in\mathbb F_p^n$. При изменении $a_j$ каждый запрещённый вектор $f$ с
$f_j\ne0$ исключает ровно одно значение
\[
a_j=-\frac{\sum_{i\ne j}a_i f_i}{f_j}\pmod p.
\]
Поэтому гистограмма запрещённых значений одновременно оценивает все $p$
вариантов координаты, сначала по числу конфликтующих проективных пар, затем по
сумме квадратов их геометрических дефицитов. Для пары координат каждое
ограничение задаёт аффинную прямую в $\mathbb F_p^2$, поэтому все $p^2$ парных
ходов также оцениваются точно за $O(p|F|)$. В режиме продолжения можно поменять
приоритеты и сначала минимизировать суммарный квадрат дефицитов, сохраняя
несколько мелких конфликтов вместо одного глубокого. При $p=139$
однокоординатный шаг уменьшил число конфликтующих пар с десяти до одной;
повторная деформация и продолжение по активному набору впервые пересекли единицу и дали
промежуточную оценку~$139$.

Для дальнейшего спуска составной индекс раскладывается на примарные блоки.
При фиксированных остальных строках лучшая строка блока выбирается глобально
из полного проективного пространства; геометрический просмотр на шаг вперёд по парам удерживает
несколько строк одного блока и для каждой точно находит лучший ответ другого.

\paragraph{Примарный Radon-поиск.} Обычный min-conflicts меняет отдельный
коэффициент формы и видит лишь локальную окрестность. Для простого $p$ мы вместо
этого оцениваем \emph{одновременно все} строки
$a\in\mathbb P^{n-1}(\mathbb F_p)$. Пусть $h(x)$ --- гистограмма остатков ещё не
отделённых запрещённых векторов. Вектор $v$ остаётся неотделённым строкой $a$ в
точности тогда, когда $a\cdot v=0$, поэтому число \emph{остающихся} конфликтов
есть конечное преобразование Радона
\[
 {\cal R}h(a)=\sum_{a\cdot x=0}h(x)
 =\frac1p\sum_{t\in\mathbb F_p}\widehat h(ta).
\]
Один $n$-мерный БПФ даёт значения ${\cal R}h$ сразу для всех строк, и наилучший
ход --- это $\arg\min_a{\cal R}h(a)$, то есть глобальный минимум по всему
проективному пространству строк, а не результат локального шага.
Повторяя такие ходы по примарным компонентам факторгруппы, мы ищем целое подпространство характеров, а не отдельные коэффициенты. Для
связи с последующей деформацией метрики используем веса
\[
 w_\beta(v)=\exp\!\left(\beta\left(1-
 \frac{D(v)}{\diam V_0}\right)\right):
\]
глубокий конфликт штрафуется сильнее пограничного.

